\chapter{Discussion}
\label{ch:discussion}

Beating the original baselines is weaker evidence for semantic grouping than
beating normalized-string counting. The original baselines differ from semantic
entropy in several ways at once, so an advantage over them cannot be traced to
any single step. Normalized-string counting differs in one step only, which
makes it the comparison that carries information about grouping. In the
sentence-answer comparisons under LLM grading, the tested Qwen entailment model
clears that stronger comparison, whereas no additional benefit is detected with
the tested DeBERTa models. The entailment model is therefore central to
interpreting the method's performance, and this chapter returns to it in every
section.

\section{What grouping by meaning adds}
\label{sec:what-clustering-is-worth}

\subsection{Normalization gives the most consistent improvement}

The decomposition in Section~\ref{sec:ladder} separates counting strings,
normalizing them, grouping them by meaning, and weighting the resulting groups
by likelihood. Normalization gives the most consistent improvement across the
tested conditions. It adds 0.0398 AUROC on average, it is positive in all 24
configuration entries, and every one of those intervals excludes zero. The two
string scores do not use an entailment model, so the 24 entries repeat 12
distinct condition--grader contrasts, and all 12 are positive. Giving those 12
equal weight still yields $+0.0372$.

Normalization is also the only step whose
difference exceeds the seed-variation reference of 0.0245, which means it is
the only step that moves AUROC further than rerunning the same configuration
under a different sampling seed. What it does is mechanical rather than
semantic: it merges differences in capitalization, punctuation, articles, and
spacing without calling another language model.

The other three steps behave differently. Counting raw strings instead of
reading sequence likelihoods is positive in only 4 of 24 comparisons, and in 8
comparisons the interval favors sequence-likelihood entropy instead. The four
positive differences are large enough to carry the overall mean to $+0.0067$,
which is a case where a positive average does not describe most of the
comparisons behind it. The claim that counting answer strings beats reading
log-likelihoods therefore holds only after the strings are normalized.

Grouping the samples by meaning adds 0.0016 AUROC over normalized-string
counting and is positive in 14 of 24 comparisons, with seven intervals favoring
normalized-string counting. Weighting the meaning groups by likelihood adds
0.0015. Both of these sit far below the normalization step and well inside the
seed range, which explains why discrete semantic entropy performs so similarly
to the likelihood-weighted version throughout this study.

The comparison with surface entropy matters because the original baselines do
not isolate grouping by meaning. Surface entropy also outperforms those
baselines in every comparison where it was measured, so a result reported only
against the original baselines would credit grouping with an improvement that
normalization already delivers. The decomposition identifies normalization,
rather than grouping alone, as the consistent source of improvement.

\subsection{The entailment model decides whether grouping adds anything}

In the short-answer comparisons, no additional benefit over normalized-string
counting is detected with the tested DeBERTa cross-encoders under LLM grading.
None of the corresponding hierarchical intervals excludes zero. This holds
across all twelve comparisons that pair a released cross-encoder with an LLM
correctness judge, so it is not one unlucky configuration.

Substituting \texttt{Qwen2.5-72B-Instruct} as the entailment model changes that
picture. On identical sentence-length answers under the Qwen judge, discrete
semantic entropy records 0.751 with the large DeBERTa-v3 model, 0.757 with the
released-code default, and 0.823 with Qwen-72B, against a surface entropy of
0.7364 that does not move at all because it never calls an entailment model.
The resulting differences are $+0.0147$, $+0.0204$, and $+0.0867$, and only the
last has an interval that excludes zero.

The stored answers are identical in
all three cases, so the entailment model is the only factor that changed. The
spread it produces, from 0.751 to 0.823, is larger than most of the method
differences this thesis measures.

Two checks support reading this as an entailment-model effect rather than an
artifact. Re-grading the same answers with Llama-3.1-70B leaves the ranking
unchanged and the sizes close, at $+0.0872$, $+0.0223$, and $+0.0139$, which
reduces the concern that one model serving as both entailment model and
correctness judge inflates the Qwen result.

Splitting the sentence-length
difference by generator shows the same dependence from another angle. Under the
Llama judge on the default entailment model, the generator with the longest
answers, Llama-3.1-8B at 667 mean characters, has the smallest difference at
$+0.012$. On the Qwen entailment model the same generator has the largest at
$+0.132$. The generations behind those two columns are the same text.

The entailment model is therefore part of the measurement rather than a setting
in the background. The recorded grid varies it across three models, which is
enough to show that the choice moves the result and not enough to say which
choice is closest to human judgments of equivalence. No human equivalence
labels were collected, so the three models can be compared with each other but
not against an outside standard.

The pattern across all 24 comparisons carries the same message. Eleven
hierarchical intervals exclude zero. The four that favor semantic entropy all
use the Qwen entailment model, and the seven that favor surface entropy all use
token-F1 correctness labels. Neither the estimator nor its comparator changes
across those eleven results. Only the entailment model and the correctness rule
do.

\subsection{What answer length explains, and what it does not}
\label{sec:length-discussion}

Farquhar et al.~\cite[Supplementary Note~7 and Supplementary Fig.~7]{semantic_entropy_nature_2024}
report a larger advantage over naive entropy for sentence-length answers than
for short phrases. The comparison in Section~\ref{sec:answer-length-results}
shows the same direction across four prompting conditions. The
discrete-semantic advantage over naive predictive entropy rises from 0.0110 for
short phrases at 11.2 mean characters, to 0.0326 at 69.8 characters, to 0.0628
for the sentence instruction without examples at 255.2 characters.

The ordering
follows the measured answer length rather than the name of the condition, which
is visible in the sentence instruction with five examples: it produces answers
of 11.4 characters and a difference of 0.0094, in line with the other short
condition.

That agreement weakens against the stronger comparator. Measured against
surface entropy instead of naive predictive entropy, the same four differences
run from $-0.0020$ to $+0.0204$, and none of the four hierarchical intervals
excludes zero. Length therefore predicts how much semantic entropy gains over
reading sequence likelihoods, and does not establish a gain over counting
normalized strings with the default entailment model.

The long-form biography evaluation shows the largest recorded gap in the study:
0.670 with large DeBERTa and 0.708 with Qwen, against 0.497 for surface and
naive sample entropy. The reason is visible in the inputs. Nearly every sampled
paragraph is a distinct string, so every string-based group holds one sample,
the frequencies are equal, and the string entropy is constant and near chance
by construction.

Meaning-based grouping restores variation in the score
(Section~\ref{sec:longform-results}), and the pooled gaps are $+0.174$ with
large DeBERTa and $+0.212$ with Qwen. Those intervals resample records while
holding generators and seeds fixed, so they are narrower in scope than the
short-answer intervals. The task, prompts, and grading target also change
together with length, so this result does not isolate answer length.

\section{Why the evaluation setup matters}
\label{sec:apparatus}

The entailment model, sampling temperature, and correctness rule can change both
the size and the direction of a method comparison. They are also different kinds
of choice. The entailment and grading checks reuse the same stored answers, so
the generated text is held fixed and only the analysis moves. The temperature
experiment instead changes the uncertainty samples while keeping the graded
answer and its LLM correctness label fixed. Neither should be described as a
change made to one fixed generation set, because only the first leaves the
samples untouched.

\subsection{The entailment model is part of the method}

Reproducing these comparisons takes three pieces of information: the entailment
model, the prompt, and the clustering rule. The label ``bidirectional
entailment'' alone fixes none of them, and the measurements above show what
that costs. The same answers and the same estimator span 0.751 to 0.823
depending on which model reads them. A reader given only the label cannot tell
where in that range a reported number sits, and cannot tell whether a gap over
a string baseline reflects the estimator or the model chosen to support it.

This thesis records all three for every comparison, which is what makes its
entailment columns comparable to one another. What it does not record is an
outside reference for the clustering rule itself. The rule is the one in the
released code, including its order dependence, and it was kept unchanged so
that the entailment-model comparison would isolate the model, so the study
measures how much the rule's output moves and not how often it is right.

\subsection{Sampling temperature changes which detector leads}
\label{sec:diversity-vs-likelihood}

At lower temperatures the sampled answers become less varied, and the detectors
that measure variation across samples have less to measure. Discrete semantic
entropy falls from 0.757 at $T=1.0$ to 0.568 at $T=0.1$, and surface entropy
from 0.736 to 0.589. Naive predictive entropy, which reads sequence likelihoods
rather than spread across samples, falls much less, from 0.694 to 0.644, and at
$T=0.1$ it leads both diversity-based estimators. The ordering therefore
reverses inside the tested range.

The graded answers and their correctness
labels are identical at every temperature, so the reversal comes from the
uncertainty samples alone. The highest temperature tested here is the most
favorable setting for the sample-based detectors in this sweep rather than a
general optimum, and a comparison reported at one temperature does not carry
over to another.

\subsection{The correctness rule can reverse the comparison}

Identical answers and identical detector scores can produce a different method
ranking when the correctness labels change. In the sentence-length condition,
surface entropy leads under token-F1, while semantic entropy has higher point
estimates under both LLM judges, and the LLM-graded intervals include zero with
the default entailment model. The short-phrase condition shows the same
reversal in miniature. The paired difference is $-0.0133$ under token-F1,
$+0.0005$ under the Qwen judge, and $+0.0027$ under the Llama judge. All seven
hierarchical intervals that favor surface entropy use token-F1 labels.

The mechanism is not subtle. Token-F1 compares an answer with a short
reference, so it can penalize an answer that is correct but long. Its low
acceptance rate in some groups also makes AUROC estimates unstable. Excluding
those groups leaves surface entropy's token-F1 advantage largely unchanged with
the default entailment model (Section~\ref{sec:grading-results}), so the effect
is not confined to a few extreme cases. The grading rule is part of what a
reported AUROC means.

\subsection{The resampling scheme changes which differences survive}

How an interval is built matters as much as what it contains. Across the same
24 comparisons, resampling records alone leaves 18 intervals that exclude zero,
with a mean width of 0.0153. Resampling the dataset--generator units as well
leaves 11, with a mean width of 0.0467, about three times wider. Seven apparent
differences disappear.

Resampling units alone identifies the same eleven, which
places the extra variation between units rather than between questions. With
three datasets and three generators there are only nine units, so a difference
that survives the record-level scheme alone describes the questions in this
grid rather than the grid itself.

\section{Comparison with Farquhar et al.\ and Kossen et al.}
\label{sec:comparison-farquhar}
\label{sec:reproduction}

The sentence-answer results agree with Farquhar et
al.~\cite{semantic_entropy_nature_2024} on the measured original baseline
comparisons and on the similar performance of the two semantic estimators.
Discrete semantic entropy leads P(True) by 0.025 on average over six
comparisons, the semantic entropy probe by 0.051, and the accuracy probe by
0.049, and it leads all three in every comparison where each was measured.
These are the comparisons the original study makes, and they reproduce here.

The agreement is limited to the configurations tested. Farquhar et
al.~\cite{semantic_entropy_nature_2024} report that P(True) tends to improve
with model size. P(True) does not show this trend within the tested Llama
models, where the 70B model records 0.711 against 0.726 for the 8B in
\texttt{chat\_0shot} and 0.819 against 0.826 in \texttt{default\_5shot}. A
comparison within one family and across two sizes is not enough to reject a
broader trend.

The baseline contrasts also do not hold semantic entropy's
entailment model fixed, because the comparator uses the default entailment
model against P(True) and \texttt{nli-deberta-v3-large} against the probes. The
missing saved training feature limits exact reproduction of the probe fits.

Kossen et al.~\cite{kossen_semantic_entropy_probes_2024} evaluate probes both
within their training tasks and on held-out tasks, and their reported advantage
over accuracy probes concerns transfer to unseen tasks. The six probe
comparisons in Section~\ref{sec:paper-baselines} use in-distribution probes,
trained and evaluated on separate questions from the same task. They do not
test that transfer finding, so the probe results here are not evidence for or
against it.

The reimplementation uses three of Farquhar et al.'s five short-answer
datasets: TriviaQA, NQ-Open, and SVAMP. It retains Mistral 7B Instruct and
tests Llama-3.1 at 8B and 70B in place of LLaMA-2 Chat. This gives direct
dataset overlap and partial model overlap while testing a smaller grid. The
generation settings also allow the answer being graded to be longer than the
samples used to estimate uncertainty, so the graded text and the text behind
the uncertainty score are not always produced under the same limits. Following
the released code makes those choices traceable, but does not remove their
effect on the results.
